Sia $\left\{x_n\right\}_{n\in\mathbb{N}}$ una serie in $\left(\mathbb{R},\left|\cdot\right|\right)$
allora possiamo definire le seguenti caratterisitche.


\begin{definition}
    \textbf{Monotonia crescente}

    Una serie si dice monotona crescente (o strettamente crescente o crescente in senso stretto)
    se e solo se
    \[
        \exists n_0 \mid \forall n \geq n_0, x_n<x_{n+1}.
    \]

    Più in generale si dice che una serie è monotona crescente in senso lato se e solo se
    \[
        \exists n_0 \mid \forall n \geq n_0, x_n\leq x_{n+1}.
    \]

    Una serie crescente in senso lato implica dunque, logicamente, che sia crescente
    anche in senso stretto.
\end{definition}

\begin{definition}
    \textbf{Monotonia decrescente}

    Una serie si dice monotona decrescente (o strettamente decrescente o decrescente in senso stretto)
    se e solo se
    \[
        \exists n_0 \mid \forall n \geq n_0, x_n>x_{n+1}.
    \]

    Più in generale si dice che una serie è monotona decrescente in senso lato se e solo se
    \[
        \exists n_0 \mid \forall n \geq n_0, x_n\geq x_{n+1}.
    \]

    Una serie decrescente in senso lato implica dunque, logicamente, che sia decrescente
    anche in senso stretto.
\end{definition}

\begin{theorem}\label{teo:mon-e-reg}
    \textbf{Monotonia e regolarità}

    Se la serie $x_n$ è definitivamente monotona (anche in senso lato), allora la serie è anche regolare.

    In particolare avremo che per tutti gli indici $n$ per cui la serie è
    monotona crescente (se è definitivamente crescente), allora
    \[
        \exists n_0 \mid n \geq n_0 \Rightarrow
        \underset{n\rightarrow+\infty}{\lim}x_n=\sup\left\{x_n\mid n \geq n_0\right\}.
    \]
    Analogamente avremo che per tutti gli indici $n$ per cui la serie è
    monotona decrescente (se è definitivamente decrescente), allora
    \[
        \exists n_0 \mid n \geq n_0 \Rightarrow
        \underset{n\rightarrow+\infty}{\lim}x_n=\inf\left\{x_n\mid n \geq n_0\right\}.
    \]
\end{theorem}

\begin{demonstration}
    \textit{di} \teoref{teo:mon-e-reg}

    \[
        \Hp x_n \text{ è definitivamente monotona crescente (risp. decrescente)}\]\[
        \Th \underset{n\rightarrow+\infty}{\lim}x_n=\sup\left\{x_n\mid n \geq n_0\right\} \text{ (risp. $\inf$)}
    \]

    \begin{enumerate}[label=(\roman*)]
        \item Se $x_n$ è definitivamente crescente (risp. decrescente), allora
        \[
            \exists n_0\mid\forall n\geq n_0, x_n \leq x_{n+1} \text{ (risp. $\geq$)}.\]\[
        \]
        \item Sapendo che esiste, fissiamo $n_0$ come valore che soddisfi la condizione di cui sopra.
        \item La serie è convergente oppure no
        \item Se è convergente
        \begin{enumerate}[label=\arabic*.]
            \item Chiamando $p$ il punto di convergenza per convergenza abbiamo che
                \[ \forall \varepsilon>0, \exists n_\varepsilon\in\mathbb{N} \mid \forall n\geq n_\varepsilon, d\left(x_n-p\right) < \varepsilon, \]
                \[ \forall \varepsilon>0, \exists n_\varepsilon\in\mathbb{N} \mid \forall n\geq n_\varepsilon, \left|x_n - p\right| < \varepsilon, \]
                \[ \forall \varepsilon>0, \exists n_\varepsilon\in\mathbb{N} \mid \forall n\geq n_\varepsilon, -\varepsilon<x_n - p< \varepsilon, \]
                \[ \forall \varepsilon>0, \exists n_\varepsilon\in\mathbb{N} \mid \forall n\geq n_\varepsilon, p-\varepsilon <x_n< p+\varepsilon. \]
            \item Per ipotesi la serie è definitivamente crescente (risp. decrescente)
                \[ \forall n \geq n_0, x_n \leq x_{n+1} \quad \text{(risp. $x_n \geq x_{n+1}$)}.\]
            \item Mostriamo che $\forall n \geq n_0, x_n \leq p$ (risp. $p \leq x_n$)
            \item Supponiamo per assurdo che $\exists k \geq n_0 \mid x_k > p$ (risp. $p > x_k$)
            \begin{enumerate}[label=\arabic{enumii}.\arabic*.]
                \item Rielaborando la supposizione avremo che \[
                    x_k > p \quad\text{(risp. $p > x_k$)},\]\[
                    x_k - p > 0 \quad\left(\text{risp. $p - x_k > 0$}\right),\]\[
                    \frac{x_k - p}{2} > 0 \quad\left(\text{risp. $\frac{p - x_k}{2} > 0$}\right),\]\[
                    \frac{x_k - p}{2} > 0 \quad\left(\text{risp. $\frac{p - x_k}{2} > 0$}\right).\]\[
                \]
                \item Per convergenza della serie, prendendo $\varepsilon_0 := \frac{x_k - p}{2}$ (risp. $\frac{p - x_k}{2}$), se $\varepsilon_0 >0$ allora, \[
                    \exists j_0\in\mathbb{N} \mid \forall j\geq j_0, p-\varepsilon_0 <x_j< p+\varepsilon_0.\]
                \item Dunque $\forall j \geq j_0$, \[
                    x_j< p+\varepsilon_0 \quad\left(\text{risp. $x_j > p-\varepsilon_0$}\right), \]\[
                    x_j< p+\frac{x_k - p}{2} \quad\left(\text{risp. $x_j > p-\frac{p - x_k}{2}$}\right), \]\[
                    x_j< \frac{2p + x_k - p}{2} \quad\left(\text{risp. $x_j > \frac{2p - \left(p - x_k\right)}{2}$}\right), \]\[
                    x_j< \frac{p + x_k}{2} \quad\left(\text{risp. $x_j > \frac{2p - p + x_k}{2}$}\right), \]\[
                    x_j< \frac{p + x_k}{2} \quad\left(\text{risp. $x_j > \frac{p + x_k}{2}$}\right), \]\[
                \]

                \item Ricordandoci che $p < x_k$(risp. $p > x_k$) \[
                    p < x_k \left(\text{risp. $p > x_k$}\right), \]\[
                    p + x_k< 2x_k \left(\text{risp. $p + x_k> 2x_k$}\right), \]\[
                    \frac{p + x_k}{2} < \frac{2x_k}{2} \quad\left(\text{risp. $\frac{p + x_k}{2} >  \frac{2x_k}{2}$}\right), \]\[
                    \frac{p + x_k}{2} < x_k \quad\left(\text{risp. $\frac{p + x_k}{2} >  x_k$}\right), \]\[
                    x_j < \frac{p + x_k}{2} < x_k \quad\left(\text{risp. $x_j > \frac{p + x_k}{2} >  x_k$}\right), \]\[
                    x_j < x_k \quad\left(\text{risp. $x_j >  x_k$}\right). \]\[
                \]

                \item Prendiamo arbitrariamente $j > \max\left\{j_0, k\right\}$, allora \[
                    \begin{cases}
                        j > j_0\\
                        j > k
                    \end{cases}\]\[
                    \begin{cases}
                        x_j < x_k \\
                        x_j > x_k 
                    \end{cases}\quad\left(\text{risp.} \begin{cases}
                        x_j > x_k \\
                        x_j < x_k 
                    \end{cases}\right)\]\[
                    \text{Assurdo (risp. Assurdo)}
                \]
                \item Dunque $\nexists k \geq n_0 \mid x_k > p$ (risp. $p > x_k$).
                \item Perciò $\forall k \geq n_0 , x_k \leq p$ (risp. $p \leq x_k$).
                \item Ovvero $\forall n \geq n_0 , x_n \leq p$ (risp. $\geq$).
            \end{enumerate}
            \qed

            \item $\forall \varepsilon>0$ abbiamo allora che \[
                \begin{cases}
                    \forall n \geq n_0 , x_n \leq p \quad\left(\text{risp. } \geq\right)\\
                    \exists n_\varepsilon\in\mathbb{N} \mid \forall n\geq n_\varepsilon, p-\varepsilon <x_n< p+\varepsilon
                \end{cases}.
            \]

            \item Prendiamo un qualsiasi $n > \max\left\{n_0, n_\varepsilon\right\}$, allora $\forall \varepsilon>0$ \[
                \begin{cases}
                    x_n \leq p \quad\left(\text{risp. } \geq\right)\\
                    p-\varepsilon <x_n< p+\varepsilon
                \end{cases},\]\[
                p-\varepsilon <x_n \leq p< p+\varepsilon  \quad\left(\text{risp. } p-\varepsilon <p \leq x_n< p+\varepsilon\right),\]\[
                p-\varepsilon <x_n \leq p\quad\left(\text{risp. } p \leq x_n< p+\varepsilon\right). \]\[
            \]

            \item Se ne deduce quindi che \[
                \begin{cases}
                    \forall n \geq n_0 , x_n \leq p \quad\left(\text{risp. } x_n\geq p\right)\\
                    \forall \varepsilon>0, \exists n \geq n_0 \mid p-\varepsilon <x_n \leq p \quad\left(\text{risp. } p \leq x_n< p+\varepsilon\right)
                \end{cases}.
            \]

            \item Da cui \[
                \begin{cases}
                    \forall n\in\left\{n\in\mathbb{N}\mid n\geq n_0\right\}, x_n \leq p \quad\left(\text{risp. } x_n\geq p\right)\\
                    \forall \varepsilon>0, \exists n\in\left\{n\in\mathbb{N}\mid n\geq n_0\right\} \mid p-\varepsilon <x_n \leq p\quad\left(\text{risp. } p \leq x_n< p+\varepsilon\right)
                \end{cases},\]\[
                \begin{cases}
                    \forall x_i\in\left\{x_n\mid n\geq n_0\right\}, x_i \leq p \quad\left(\text{risp. } x_n\geq p\right)\\
                    \forall \varepsilon>0, \exists x_\varepsilon\in\left\{x_n\mid n\geq n_0\right\} \mid p-\varepsilon <x_\varepsilon \leq p\quad\left(\text{risp. } p \leq x_n< p+\varepsilon\right)
                \end{cases}.
            \]

            \item Di conseguenza per \defref{def:sup-lim-rea} e \defref{def:inf-lim-rea} segue necessariamente che$p = \sup\left\{x_n\mid n\geq n_0\right\}$ (risp. $\inf$).
        \end{enumerate}

        \item Se non è convergente
        \begin{enumerate}[label=\arabic*.]
            \item Mostriamo separatamente che una successione definitivamente crescente (risp. decrescente) e non convergente ha come limite $+\infty$ (risp. $-\infty$) e che l'insieme
            degli elementi per cui è crescente (risp. decrescente) non è superiormente (risp. inferiormente) limitato, ovvero il suo $\sup$ (risp. $\inf$) è $+\infty$ (risp. $-\infty$).

            \item Iniziamo mostrando che la successione così definita risulta divergente a $+\infty$ (risp. $-\infty$).
        \end{enumerate}
    \end{enumerate}
    \qed
\end{demonstration}

\begin{theorem}\label{teo:com-rea}
    \textbf{Completezza di $\mathbb{R}$}

    $\left(\mathbb{R}, \left|\cdot\right|\right)$ è completo, ovvero tutte le successioni
    di Cauchy in $\left(\mathbb{R}, \left|\cdot\right|\right)$ convergono a un elemento di $\mathbb{R}$.
\end{theorem}

\begin{demonstration}
    \textit{di} \teoref{teo:com-rea}

    [dimostrazione]
\end{demonstration}